\documentclass[11pt]{article}
\usepackage[T1]{fontenc}
\usepackage{amsmath,amssymb,amsthm}
\usepackage[margin=1in]{geometry}
\usepackage{xurl}
\usepackage[hidelinks]{hyperref}
\newtheorem{theorem}{Theorem}
\newtheorem{lemma}[theorem]{Lemma}
\newtheorem{proposition}[theorem]{Proposition}
\newcommand{\Spec}{\operatorname{Spec}}
\newcommand{\m}{\mathfrak m}
\title{Direct algebraic reductions for the \'etale argument\\
\large What can replace the dynamical rereading?}
\author{}
\date{31 October 2026}
\begin{document}
\maketitle

The direct flatness implication in this note is now a checked Unit
theorem in \path{sag/etale-direct}. The first part gives mathematical
arguments; the last part gives their formalisation correspondence.
The residual-cover construction remains a proposal. The two targets must be distinguished: unique square-zero
lifting, and a finite cover by \emph{one-variable monic} standard charts.

\part{Mathematics}

\section{The lesson from the direct Horrocks argument}

The useful simplification in the division proof of rank-one Horrocks
was to replace residue-field reasoning with identities over the original
ring. An annihilator $\mu\equiv1\bmod J$ kills a finite module. The two
explicit opens at $\mu$ and $1-\mu$ give either the answer or a strict
decrease in the polynomial degree. Thus the branching is bounded by an
algebraic measure, rather than being obtained by interpreting an arbitrary
proof about a generic prime.

The analogous questions here are: which uses of a prime merely select
among finitely many available certificates; which residue-field arguments
can be replaced by equations over the original ring; and where would a
termination measure for a direct covering argument come from?

\section{Removing the selected prime from the residual construction}

Suppose first that the base is residually discrete local, with residue
field $k$, and that the completed field theorem supplies
\[
       C=E/\m E\simeq \prod_i k[T]/(P_i),
\]
with $P_i$ monic and separable. Zero factors may be retained.
Write $e_i$ for the product idempotents and $x_i$ for generators of the
corners, whose identity elements are $e_i$.

Instead of selecting the factor outside a prime and then selecting $x_i$
or $x_i-1$, retain both elements
\[
         y_{i,0}=e_i x_i,\qquad y_{i,1}=e_i(x_i-1).
\]
Here $x_i$ may be any lift to $C$ of the generator of the corner. Then
\[
          y_{i,0}-y_{i,1}=e_i,
          \qquad \sum_i(y_{i,0}-y_{i,1})=1.
\]
This is an actual finite cover certificate. No prime membership test
is needed. The translated generator has translated monic separable
polynomial, so it carries the same kind of input to the lifting argument.

The useful invariant for each subsequent denominator $h_{i,j}$ is not
``outside the selected prime'', but an explicit relation
\[
               y_{i,j}^{N_{i,j}}\in (\overline h_{i,j})\subseteq C.
\]
For example, a determinant congruence $p\equiv y^r\bmod\m$ supplies
such information directly. An inverse to a derivative after inverting
$y$ supplies it after clearing a power of $y$. These calculations replace
uses of primality with polynomial and localisation identities.

\begin{lemma}\label{lem:powers}
If finitely many $y_j$ generate the unit ideal and $y_j^{N_j}\in(h_j)$
with $N_j\geq1$, then the $h_j$ generate the unit ideal.
\end{lemma}
\begin{proof}
Raise an identity $1=\sum_j a_j y_j$ to the power
$1+\sum_j(N_j-1)$. Every resulting monomial contains some $y_j^{N_j}$,
and substitute its given expression as a multiple of $h_j$.
The empty-family case is the zero ring.
\end{proof}

This suggests a precise replacement for the selected-prime local
interface: construct finitely many standard charts $E[1/h_j]$, together
with an identity
\begin{equation}\label{eq:residual-cover}
             1=\sum_j a_jh_j+\sum_\ell m_\ell b_\ell,
             \qquad m_\ell\in\m.
\end{equation}
The proposal is to obtain this from the existing product theorem and
the existing candidate construction, using the power invariant above.
This complete residual-cover adapter is not yet formalised. The
first-unit and polynomial degree-lowering algorithms still use the
discrete residue field; retaining all factors does not remove that input.

\section{Why a residual cover is not a global cover}

Equation~\eqref{eq:residual-cover} gives $h=\sum_j a_jh_j\in1+\m E$.
The charts cover $E[1/h]$. It does not follow that $h$ is a unit in $E$.

For example, take $R=k[t]_{(t)}$ and
\[
              E=R\times R[1/t],\qquad \m=(t).
\]
Both factors are \'etale over $R$, but $E/\m E\simeq k$. The idempotent
$h=(1,0)$ satisfies
\[
              1-h=(0,1)=t(0,1/t)\in\m E.
\]
Nevertheless $h$ is not a unit: its open misses the entire second factor.
Thus even the full product decomposition of the closed fibre does not
give a global cover of $E$.

Equation~\eqref{eq:residual-cover} does give a cover by the $h_j$ together
with the images of the finitely many $m_\ell$. Continuing on the latter
opens changes the base to $R[1/m_\ell]$. There is presently no degree or
rank bound established for recursively repeating this procedure. An
unjustified iteration here would simply hide the remaining local--global
problem. For finite $R$-algebras, the Jacobson-radical property of $\m E$
would remove this particular obstruction; the original $E$ need not be
finite over $R$.

\section{A direct Jacobian source over any base}

There is a stronger simplification for the implication from flatness and
unramifiedness to square-zero lifting. It needs neither residue fields
nor a cover.

Let
\[
       E=R[X_1,\ldots,X_n]/(f_1,\ldots,f_m),
       \qquad x_i=X_i\bmod(f_1,\ldots,f_m),
\]
and assume formal unramifiedness. The already available diagonal
certificate gives $\epsilon\in E\otimes_R E$ with
\[
       \mu(\epsilon)=1,\qquad
       \epsilon(1\otimes x_i-x_i\otimes1)=0.
\]
Idempotence of $\epsilon$ is not needed in the next calculation.

\begin{lemma}[Jacobian certificate from the diagonal]\label{lem:jac}
Put $J_{jk}=(\partial f_j/\partial X_k)(x)$. There is a matrix
$C=(c_{ij})\in M_{n\times m}(E)$ with $CJ=I_n$.
\end{lemma}
\begin{proof}
Regard the tensor algebra as
\[
       E[Y_1,\ldots,Y_n]/(f_1(Y),\ldots,f_m(Y)),
\]
and choose a polynomial $p(Y)$ representing $\epsilon$. Then $p(x)=1$.
For each $i$, kernel membership supplies polynomials $a_{ij}(Y)$ such
that, in the polynomial ring over $E$,
\[
       p(Y)(Y_i-x_i)=\sum_j a_{ij}(Y)f_j(Y).
\]
Differentiate with respect to $Y_k$ and evaluate at $Y=x$. Terms
containing $Y_i-x_i$ or $f_j(x)$ vanish, giving
\[
              \delta_{ik}=\sum_j a_{ij}(x)J_{jk}.
\]
Take $c_{ij}=a_{ij}(x)$. Only finitely many representatives and kernel
witnesses are used.
\end{proof}

Lift the finitely many $c_{ij}$ to polynomials $\widetilde c_{ij}\in R[X]$
and put
\[
       F_i=\sum_j\widetilde c_{ij}f_j,\qquad
       \Delta=\det\left(\frac{\partial F_i}{\partial X_k}\right).
\]
Because every $f_j$ vanishes in $E$, the Jacobian of the $F_i$ becomes
$CJ=I_n$ there. In particular $\Delta(x)=1$.

\begin{proposition}\label{prop:source}
There is a surjective $R$-algebra map
\[
       u:S=\bigl(R[X_1,\ldots,X_n]/(F_1,\ldots,F_n)\bigr)[1/\Delta]
           \longrightarrow E.
\]
The algebra $S$ is finitely presented and formally \'etale.
\end{proposition}
\begin{proof}
The map sends $X_i$ to $x_i$ and $1/\Delta$ to $1$. It is surjective
because the $x_i$ generate $E$.

For formal \'etaleness, consider an $R$-algebra map $S\to D/N$, with
$N^2=0$. Lift the finitely many coordinates to a tuple $a\in D^n$.
The determinant of $J_F(a)$ is a unit modulo $N$, hence a unit in $D$.
Put
\[
                    b=a-J_F(a)^{-1}F(a).
\]
The correction belongs to $N^n$. The exact first-order Taylor identity
modulo $N^2$ gives $F(b)=0$. Also $\Delta(b)$ is a unit because its
image modulo $N$ is unchanged. This defines a lift. If $b'$ is another
lift, Taylor's identity gives $J_F(b)(b'-b)=0$, hence $b'=b$.
For $n=0$, take the empty determinant $1$ and $S=R$.
\end{proof}

This construction avoids a cover by Jacobian minors: suitable linear
combinations of the original relations already give an invertible square
Jacobian everywhere on the image of $E$. The statement is a multivariable
Jacobian presentation. It does not assert a monic one-variable standard
presentation. Compare the distinction in Stacks, Sections
\href{https://stacks.math.columbia.edu/tag/00U0}{10.143} and
\href{https://stacks.math.columbia.edu/tag/0G1A}{10.144}.

\section{The short ending using flatness}

Use the already proved flatness transfer principle: if $S/R$ is finitely
presented and formally unramified, then any $R$-flat $R$-algebra $E$
receiving a map from $S$ is flat over $S$.

\begin{lemma}[Lifting through a flat quotient]\label{lem:quotient}
Let $u:S\twoheadrightarrow E$ be a surjective map of $R$-algebras.
Suppose $S/R$ is formally smooth and $E$ is flat over $S$. Then $E/R$
is formally smooth.
\end{lemma}
\begin{proof}
Write $K=\ker(u)$. A flat quotient has the following explicit property:
for each $z\in K$ there is $b\in K$ such that $z=zb$.
Given a square-zero ideal $N\subseteq D$ and a map $\alpha:E\to D/N$,
lift $\alpha u$ to $v:S\to D$ by formal smoothness of $S$.
For $z\in K$, both $v(z)$ and $v(b)$ belong to $N$, and therefore
\[
                   v(z)=v(z)v(b)=0.
\]
Thus $v$ factors through $E$, providing a lift of $\alpha$.
No finite generation of $K$, chosen idempotent generator, or
determinant trick is needed at this step.
\end{proof}

\begin{theorem}\label{thm:flat}
A finitely presented, flat, formally unramified $R$-algebra is formally
\'etale, over an arbitrary commutative ring $R$.
\end{theorem}
\begin{proof}
Construct $S\twoheadrightarrow E$ by Proposition~\ref{prop:source}.
Flatness transfer gives $S$-flatness of $E$. Lemma~\ref{lem:quotient}
gives existence of square-zero lifts; unramifiedness gives uniqueness.
\end{proof}

This proves the flatness implication directly, without ZMT, residue
fields, primitive elements, prime ideals, or local--global covering.
It does not prove that the multivariable algebra $S$, or its quotient
$E$, has a finite cover by monic one-variable standard algebras.
That is the stronger conclusion still under investigation.

\part{Connection with Unit}

\section{The checked arbitrary-base theorem}

The whole direct proof is formalised in \path{sag/etale-direct}, including
the diagonal-to-Jacobian step, the constructed source, and the lifting
endpoint. In \texttt{FlatUnramifiedEtale.Over[R,setR]}, the theorem
\path{Presented[T,presentation,unramified,flat].theorem} has precisely
these hypotheses: an actual finite presentation of $T$, formal
unramifiedness, and $R$-flatness. Its conclusion is the existing
\path{FormallyEtale.Etale T}.

With a merely given finite presentation, \path{At[...].formally-etale}
still returns an actual unique lifting function. The endpoint
\path{At[...].etale} keeps the presentation-containing conclusion under
truncation. Uniqueness makes each type of lifts a proposition, permitting
elimination of the finite representative choices without any choice axiom.

\section{Correspondence with the mathematical proof}

\begin{enumerate}
\item \path{FlatQuotientLift.seg} proves Lemma~\ref{lem:quotient}, reusing
\path{FlatQuotient} and \path{EtaleCase.Surjection.Factor}. No finite
kernel generators or idempotent generator are introduced.
\item \path{ExpressionJets.seg} proves directional differentiation,
linearity and the first-order Taylor identity for expression syntax.
\path{FreePolynomialJets.seg} defines differentiation on an actual free
polynomial algebra by evaluating into the constructed dual-number ring.
\path{MappedJets.seg} proves compatibility with coefficient change.
\item \path{DiagonalJacobian.seg} proves Lemma~\ref{lem:jac}. It constructs
$E[Y]$ and its quotient by the original mapped equations, then maps the
existing diagonal tensor certificate into this quotient. The effective
quotient gives the finite coefficient witnesses directly. Thus it does
not require a separate ring isomorphism identifying two presentations of
the tensor product. \path{PresentationKernel} is used for equation-indexing
helpers; its main kernel theorem is unnecessary here.
\item \path{JacobianTools.seg} supplies gradient and determinant
naturality and linear solving by the existing right adjugate.
\path{JacobianLifting.seg} constructs the finite square-system algebra
and its localisation, and proves unique square-zero lifting. This works
for every number of variables, including zero.
\item \path{UnramifiedEtaleSource.seg} lifts the finite coefficients of
$C$, constructs the $F_i$, proves that their Jacobian becomes $I$ in $E$,
and constructs the surjection $S\twoheadrightarrow E$. Its source
certificate exists merely; no flatness is used at this stage.
\item \path{FlatUnramifiedEtale.seg} composes the source theorem with
\path{FlatDiagonalTransfer} and \path{FlatQuotientLift}, obtaining
Theorem~\ref{thm:flat} over the original arbitrary base.
\end{enumerate}

The source and all the representative choices remain inside truncations.
The new development introduces no axioms. It has only the existing
univalence, replacement and surjection parameters; there is no ZMT,
field structure, local-ring, prime-ideal or decidability parameter.
All reused proof files are relative links and no shared proof or checker
file was modified.

\section{Validation}

A fresh checker run checks $5834$ declarations with no failures in
$2.26$ seconds. The eleven owned proof modules contain $557$ declaration
events, including nineteen open directives. Each new declaration takes
at most $0.02$ seconds and opens at most $32$ distinct definitions.
The frozen opening specification covers all $538$ named new declarations;
every entry reports no additional openings. This includes concrete
quotient carrier and operation aliases.

Regressions cover the empty presentation, $X=a$ and redundant equations,
and arbitrary principal localisations. Concrete instances use
$\mathbb Z/4\mathbb Z$, including localisation at its nilpotent element
$2$. A nonsymmetric two-variable Jacobian checks the orientation and the
right-adjugate solver. Mutations of the Newton sign, the defining linear
combinations and the Jacobian orientation are rejected in isolated copies.

From \path{unit/}, run \texttt{./segc sag/etale-direct all}.
The directory's \path{README.txt} gives the exact exports and audit commands.

\section{Assessment for the original global structure theorem}

There are now three distinct tasks rather than a single opaque
dynamical transformation:
\begin{enumerate}
\item The direct flatness implication above is now checked over an arbitrary
base, with its source and all intermediate certificates constructed.
\item The prime-free residual-cover adapter removes factor selection and
the $x$ versus $x-1$ decision from the local structure proof. It reuses
the completed field product theorem, but still has a discrete residue
field and only produces a cover modulo $\m E$.
\item The full standard-cover theorem still needs a direct finite
one-variable chart construction over an arbitrary base, or a justified
finite local--global argument. Neither the Jacobian presentation nor
the residual-cover identity supplies that step. The two chart descent
lemmas from the original note remain independently useful.
\end{enumerate}

The analysis therefore supports simplifying the algebra before building
a general dynamical framework. It gives a definite shortcut for the
flatness implication and a sharper, genuinely smaller target for the
remaining global structure problem; it does not claim that all
dynamical reasoning has been eliminated from that stronger problem.

\end{document}
